\chapter{Resolver problemas}\label{ch:g3:problems}

Una cosa es conocer las cuatro operaciones y otra saber \emph{cuál} pide
cada historia. Este capítulo da un método para los problemas con
enunciado --- los mismos cuatro pasos siempre --- y enseña a
leer información en tablas y pictogramas.

\section{Los cuatro pasos}

\begin{method}[Resolver un problema con enunciado]\label{met:g3:problems:steps}
\begin{enumerate}
  \item \emph{Lee} el problema dos veces; di con tus palabras qué se
        pregunta y qué se sabe.
  \item \emph{Elige} la operación: haz un esquema rápido o un dibujo de
        barras si la elección no es evidente.
  \item \emph{Calcula}, en columna o de cabeza, con una estimación como
        red de seguridad.
  \item \emph{Responde} con una frase completa que incluya la unidad ---
        y comprueba que la respuesta tiene sentido (un niño no pesa
        $300$ kg).
\end{enumerate}
\end{method}

\begin{example}[¿Qué operación?]\label{ex:g3:problems:which}
Las palabras dan pistas sobre la operación y el esquema decide:
\begin{itemize}
  \item \emph{juntar / en total}: \omterm{def:g1:addition:def}{suma};
  \item \emph{quitar / cuántos más / qué falta}:
        resta;
  \item \emph{tantos grupos de tantos / tantas veces}:
        multiplicación;
  \item \emph{repartir a partes iguales / cuántos grupos}: división.
\end{itemize}
\end{example}

\begin{omfigure}
\begin{tikzpicture}[scale=0.62]
  % barra de la suma
  \draw[thick, fill=omDef!20] (0,0) rectangle (3.4,0.8);
  \draw[thick, fill=omThm!20] (3.4,0) rectangle (5.8,0.8);
  \node[font=\small] at (1.7,0.4) {$26$};
  \node[font=\small] at (4.6,0.4) {$17$};
  \draw[Stealth-Stealth] (0,-0.35) -- (5.8,-0.35);
  \node[below, font=\small] at (2.9,-0.4) {¿cuántos en total?: $26 + 17$};
  % barra de comparacion
  \begin{scope}[xshift=9.4cm]
  \draw[thick, fill=omDef!20] (0,0.9) rectangle (5.2,1.7);
  \node[font=\small] at (2.6,1.3) {$52$};
  \draw[thick, fill=omThm!20] (0,0) rectangle (3.5,0.8);
  \node[font=\small] at (1.75,0.4) {$35$};
  \draw[Stealth-Stealth] (3.5,0.4) -- (5.2,0.4);
  \node[below, font=\small] at (2.6,-0.15)
    {¿cuántos más?: $52 - 35$};
  \end{scope}
\end{tikzpicture}

{\small Los dibujos de barras hacen visible la operación: partes una junto a otra
piden una \omterm{def:g1:addition:def}{suma}; dos barras comparadas piden una resta.}
\end{omfigure}

\begin{example}[Un problema de dos pasos]\label{ex:g3:problems:twostep}
«Maya compra $3$ paquetes de $6$ botellas y se bebe $2$ botellas. ¿Cuántas
botellas quedan?»
\begin{enumerate}
  \item qué se pregunta: el número de botellas que quedan. Se sabe: $3$
        paquetes de $6$ y $2$ bebidas;
  \item esquema: $3$ grupos de $6$ y luego quitar $2$. ¡Dos
        operaciones!
  \item cálculo: $3 \times 6 = 18$ y después $18 - 2 = 16$;
  \item frase: quedan $16$ botellas.
\end{enumerate}
No te pares nunca después de la primera operación: relee la pregunta para ver
si la historia ha terminado.
\end{example}

\begin{example}[Un número que sobra]\label{ex:g3:problems:trap}
«Un panadero, que tiene $52$ años, hornea $120$ panecillos y vende $85$.
¿Cuántos panecillos le quedan?» La edad está ahí para distraerte: la
respuesta solo necesita $120 - 85 = 35$ panecillos. Antes de calcular, ordena
siempre los números: ¿cuáles usa de verdad la pregunta?
\end{example}

\section{Leer tablas y pictogramas}

\begin{example}[Una tabla]\label{ex:g3:problems:table}
Libros prestados en la biblioteca:
\begin{center}
\renewcommand{\arraystretch}{1.3}
\begin{tabular}{l|cccc}
 & cómics & novelas & ciencia & total \\
\hline
semana 1 & $14$ & $9$ & $6$ & $29$ \\
semana 2 & $11$ & $12$ & $8$ & $31$ \\
\end{tabular}
\end{center}
Lectura: semana 2, novelas: $12$. Cálculo a partir de la tabla: entre las dos
semanas se prestaron $14 + 11 = 25$ cómics; la semana con más préstamos
fue la 2 ($31 > 29$).
\end{example}

\begin{example}[Un pictograma]\label{ex:g3:problems:pictogram}
En un pictograma, un símbolo representa una cantidad fija --- ¡lee
primero la clave!

\begin{omfigure}
\begin{tikzpicture}[scale=0.62]
  \node[left, font=\small] at (0,2.4) {lunes};
  \foreach \i in {0,1,2} \fill[omDef] (0.5+\i*0.8,2.4) circle (0.26);
  \node[left, font=\small] at (0,1.2) {martes};
  \foreach \i in {0,1,2,3,4} \fill[omDef] (0.5+\i*0.8,1.2) circle (0.26);
  \node[left, font=\small] at (0,0) {miércoles};
  \foreach \i in {0,1} \fill[omDef] (0.5+\i*0.8,0) circle (0.26);
  \fill[omDef] (0.5+1.6,0) circle (0.26);
  \node[right, font=\small] at (5.4,1.2)
    {clave: un punto $= 4$ pasteles vendidos};
\end{tikzpicture}

{\small Un pictograma: el lunes hay $3$ puntos, o sea, $3 \times 4 = 12$
pasteles; el martes $5$ puntos, $20$ pasteles; el miércoles $3$ puntos, $12$
pasteles.}
\end{omfigure}
\end{example}

\section{Ejercicios}

\begin{exercise}[$\star$]\label{exo:g3:problems:1}
Para cada historia, \emph{di} solamente la operación (sin calcular):
un rebaño de $48$ ovejas se junta con otro de $37$; \ $56$ cerezas repartidas
entre $7$ niños; \ $9$ cajas de $8$ lápices; \ Ali tenía $71$ cromos
y regaló $18$.
\end{exercise}

\begin{exercise}[$\star$]\label{exo:g3:problems:2}
Resuelve con los cuatro pasos: «Un aparcamiento tiene $246$ coches por la
mañana y a mediodía llegan $158$ más. ¿Cuántos coches hay ahora?»
\end{exercise}

\begin{exercise}[$\star$]\label{exo:g3:problems:3}
Resuelve: «Una cuerda de $90$ m se corta en trozos de $9$ m. ¿Cuántos
trozos salen?»
\end{exercise}

\begin{exercise}[$\star$]\label{exo:g3:problems:4}
Resuelve: «Una entrada cuesta $8$. ¿Cuánto cuestan $26$ entradas?»
\end{exercise}

\begin{exercise}[$\star$]\label{exo:g3:problems:5}
Dibuja un esquema de barras y resuelve: «Elena tiene $64$ pegatinas y Hugo tiene
$29$ menos. ¿Cuántas tiene Hugo?»
\end{exercise}

\begin{exercise}[$\star$]\label{exo:g3:problems:6}
Dos pasos: «Una caja lleva $45$ manzanas. Un frutero recibe $6$ cajas
y vende $178$ manzanas. ¿Cuántas manzanas le quedan?»
\end{exercise}

\begin{exercise}[$\star$]\label{exo:g3:problems:7}
Con la tabla del \cref{ex:g3:problems:table}: ¿cuántos libros de ciencia
se prestaron entre las dos semanas? ¿Cuántas novelas más
en la semana 2 que en la 1?
\end{exercise}

\begin{exercise}[$\star$]\label{exo:g3:problems:8}
Con el pictograma del \cref{ex:g3:problems:pictogram}: ¿cuántos
pasteles se vendieron entre los tres días? ¿Qué día se vendieron
más pasteles?
\end{exercise}

\begin{exercise}[$\star$]\label{exo:g3:problems:9}
Este problema tiene un número que sobra --- búscalo y después resuelve: «Un
libro de $210$ páginas cuesta $12$. Leo lee $35$ páginas al día. ¿Cuántos
días necesita para terminar el libro?»
\end{exercise}

\begin{exercise}[$\star\star$]\label{exo:g3:problems:10}
«Una familia de $2$ adultos y $3$ niños visita el museo. La entrada de adulto
cuesta $9$ y la de niño $4$. Pagan con un billete de $50$.
¿Cuánto les devuelven?» (Tres pasos: adultos, niños y después
el cambio.)
\end{exercise}

\begin{exercise}[$\star\star$]\label{exo:g3:problems:11}
Inventa tu propio problema de dos pasos cuya respuesta sea $100$, escribe su
solución con los cuatro pasos y pruébalo con un compañero.
\end{exercise}
